\begin{exercise}[Forward price with deterministic dividends]
\emph{Restatement.}  Find the no-arbitrage forward delivery price for an
asset \(S\) maturing at \(T\) when the asset pays deterministic cash dividends
before maturity.
\end{exercise}

\begin{solution}
Assume first that there is one deterministic dividend \(D>0\) paid at
\(\bar t\in(0,T)\).  A long forward entered at time \(0\) receives the asset at
time \(T\), but does not receive the dividend paid before \(T\).  Therefore the
correct comparison is the \emph{prepaid forward}: buy the asset today and
finance the dividend by borrowing its present value \(De^{-r\bar t}\).  When
the dividend is paid at \(\bar t\), it repays exactly that loan.

Thus the time-0 prepaid forward value is
\[
  S_0-De^{-r\bar t}.
\]
If the ordinary forward delivery price is \(F_0\), its time-0 value is zero
when the present value of paying \(F_0\) at \(T\) equals the prepaid forward
value:
\[
  F_0e^{-rT}=S_0-De^{-r\bar t}.
\]
Hence
\[
  \boxed{
  F_0=(S_0-De^{-r\bar t})e^{rT}
      =S_0e^{rT}-De^{r(T-\bar t)}.
  }
\]
If the delivery price were larger, one could sell the forward and buy the
prepaid forward replication; if it were smaller, one could buy the forward
and short the replication.  This is the standard AOA comparison principle of
\polyref{Cor.~1.1}.

For deterministic dividends \(D_1,\ldots,D_m\) paid at
\[
  0<\bar t_1<\cdots<\bar t_m<T,
\]
the present value of the dividends is additive.  The time-0 forward delivery
price is therefore
\[
  \boxed{
  F_0
  =
  \left(S_0-\sum_{j=1}^m D_j e^{-r\bar t_j}\right)e^{rT}
  =
  S_0e^{rT}
  -
  \sum_{j=1}^m D_j e^{r(T-\bar t_j)}.
  }
\]
More generally, at an intermediate date \(t<T\), only future dividends matter:
\[
  F_t
  =
  \left(
    S_t-\sum_{\bar t_j>t}D_j e^{-r(\bar t_j-t)}
  \right)e^{r(T-t)}.
\]
\end{solution}
